\documentclass[11pt,a4paper]{article}
\usepackage[utf8]{inputenc}
\usepackage[T1]{fontenc}
\usepackage[spanish]{babel}
\usepackage[top=2cm,bottom=2cm,left=2.2cm,right=2.2cm]{geometry}
\usepackage{booktabs}
\usepackage{array}
\usepackage{xcolor}
\usepackage{hyperref}
\usepackage{fancyhdr}
\usepackage{titlesec}
\usepackage{enumitem}
\usepackage{mdframed}
\usepackage{tabularx}
\usepackage{parskip}

\definecolor{primaryblue}{RGB}{33,150,243}
\definecolor{criticalred}{RGB}{244,67,54}
\definecolor{normalgreen}{RGB}{76,175,80}
\definecolor{warningorange}{RGB}{255,152,0}
\definecolor{darkgray}{RGB}{55,55,55}

\setlength{\parskip}{4pt}
\setlength{\parindent}{0pt}

\pagestyle{fancy}
\fancyhf{}
\rhead{\small EdgeLeak AI}
\lhead{\small Informe Ejecutivo -- Estado del Arte}
\cfoot{\small\thepage}
\renewcommand{\headrulewidth}{0.4pt}

\hypersetup{colorlinks=true,linkcolor=primaryblue,urlcolor=primaryblue,citecolor=primaryblue}

\titleformat{\section}{\large\bfseries\color{primaryblue}}{}{0em}{\thesection.\enspace}[\vspace{-2pt}\textcolor{primaryblue}{\rule{\linewidth}{0.6pt}}]
\titleformat{\subsection}{\normalsize\bfseries\color{darkgray}}{}{0em}{\thesubsection.\enspace}
\titlespacing{\section}{0pt}{8pt}{4pt}
\titlespacing{\subsection}{0pt}{5pt}{2pt}

% =====================================================================
\begin{document}

\begin{center}
  {\Large\textbf{Informe Ejecutivo: Estado del Arte}}\\[2pt]
  {\large\textcolor{primaryblue}{\textbf{EdgeLeak AI}} — Sistema Inteligente de Detección
  Temprana de Microfugas Domésticas}\\[2pt]
  {\small Revisión de 13 referentes · Brechas identificadas · Posicionamiento de la propuesta}
\end{center}
\vspace{-4pt}
\textcolor{primaryblue}{\rule{\linewidth}{1.2pt}}
\vspace{2pt}

% =====================================================================
\section{Universo de Referentes Analizados}

Se revisaron \textbf{13 referentes} (papers IEEE/MDPI, tesis, prototipos open source y
documentación técnica) publicados entre 2016 y 2025. Ninguno aborda el nicho exacto del
proyecto: \textit{detección predictiva de microfugas en acoples domésticos bajo el
fregadero, con IA en el dispositivo y costo inferior a \$20~USD}.

\vspace{4pt}
\begin{tabularx}{\linewidth}{@{} >{\small}l >{\small}X >{\small}X @{}}
\toprule
\textbf{\#} & \textbf{Referente} & \textbf{Limitación central} \\
\midrule
1 & Anónimo (s.f.) — Prototipo Arduino doméstico & Reactivo: alarma solo cuando hay charco visible \\
2 & De Lima Rosado (2025) — Seguridad IoT & Solo ciberseguridad; no aborda detección física \\
3 & Islam et al. (2023) — IoT + ML, IEEE & IA en nube; sin funcionamiento offline \\
4 & Kang et al. (2023) — Vibración MDPI & Entorno industrial controlado; falsos positivos en cocina \\
5 & Edge Impulse (2023-24) — TinyML ESP32 & Propósito general (motores); no calibrado para fluidos \\
6 & Escobar \& López (2016) — Sensor YF-S201 & Turbina ciega a caudales $<$0.5 L/min; sin IA \\
7 & Rubiños et al. (2023) — Smart Cities & Escala macro-acueducto; servidores de alto cómputo \\
8 & Farah \& Shahrour (2024) — Review MDPI & 100\% redes municipales; vacío total en Smart Home \\
9 & Corado (2025) — Smart Cities UOC & IA en servidores externos; no aplicable al hogar \\
10 & Islam et al. (2022) — Review IEEE & Tuberías subterráneas; ignora ruido de cocina \\
11 & Sagat (s.f.) — ESP8266 GitHub & Sensor pasivo reactivo; sin IA \\
12 & Dixit et al. (2017) — Presión/Nivel IRJET & Insensible a goteos lentos donde la presión no cae \\
13 & Bruno et al. (2021) — Datalogging MDPI & Registro diferido en SD; sin alerta en tiempo real \\
\bottomrule
\end{tabularx}

% =====================================================================
\section{Brechas Identificadas en el Estado del Arte}

\subsection{Brecha 1 — Lo doméstico es reactivo; lo predictivo es solo macro}

\begin{mdframed}[linecolor=criticalred,linewidth=0.8pt,backgroundcolor=red!3,
  innertopmargin=4pt,innerbottommargin=4pt,innerrightmargin=6pt,innerleftmargin=6pt]
\textbf{¿Qué no se resuelve?} No existe ninguna solución que lleve la detección predictiva
de fugas al nicho residencial de bajo costo.
\end{mdframed}

Los sensores domésticos comerciales (Xiaomi, Tuya, Sagat~[11], Anónimo~[1]) operan con
lógica binaria: detectan agua acumulada cuando el daño ya ocurrió.
Las soluciones predictivas con IA (Rubiños~[7], Farah~[8], Corado~[9], Islam~[10])
están diseñadas para acueductos municipales, requieren sensores industriales y servidores
de cómputo. \textbf{La franja Smart Home predictiva de bajo costo no tiene ocupante.}

\subsection{Brecha 2 — Acústica sin contexto genera falsas alarmas}

\begin{mdframed}[linecolor=warningorange,linewidth=0.8pt,backgroundcolor=orange!3,
  innertopmargin=4pt,innerbottommargin=4pt,innerrightmargin=6pt,innerleftmargin=6pt]
\textbf{¿Qué no se resuelve?} Ningún referente combina señal acústica con señal de
flujo para eliminar falsos positivos en el entorno doméstico.
\end{mdframed}

Kang~[4] e Islam~[3] confirman que los sensores acústicos/piezoeléctricos son el método
de mayor precisión para fugas ocultas. Sin embargo, sus pruebas son en entornos
industriales sin ruido cotidiano. En una cocina, lavar platos genera la misma firma
vibratoria que un goteo. \textbf{Ningún referente implementa Sensor Fusion
(vibración + flujo) como interruptor de contexto} para saber si el ruido es uso normal
o fuga real.

\subsection{Brecha 3 — La IA requiere la nube, exponiendo privacidad y generando latencia}

\begin{mdframed}[linecolor=primaryblue,linewidth=0.8pt,backgroundcolor=blue!3,
  innertopmargin=4pt,innerbottommargin=4pt,innerrightmargin=6pt,innerleftmargin=6pt]
\textbf{¿Qué no se resuelve?} Ningún referente ejecuta inferencia de IA localmente en
el microcontrolador para un sistema de detección de fugas domésticas.
\end{mdframed}

Islam~[3], Corado~[9] y Rubiños~[7] dependen de servidores externos para correr sus
modelos. Esto implica: (a)~el sistema queda ciego si falla el router, (b)~el audio del
hogar viaja a servidores de terceros (privacidad), (c)~la latencia de respuesta es alta.
Bruno~[13] va al extremo opuesto: registra en SD para análisis diferido, sin ninguna
alerta en tiempo real. \textbf{El paradigma Edge AI —inferencia en el dispositivo— no
ha sido aplicado a este dominio específico.}

% =====================================================================
\section{Posicionamiento de EdgeLeak AI}

\vspace{2pt}
\begin{tabularx}{\linewidth}{@{} >{\small\bfseries}p{3.8cm} >{\small}X @{}}
\toprule
\textbf{Brecha del estado del arte} & \textbf{Respuesta de EdgeLeak AI} \\
\midrule
Preventivo-doméstico inexistente &
  Modelo TinyML entrenado con firmas acústicas de goteo real, desplegado en ESP32
  (\textasciitilde\$5 USD). Detecta la turbulencia del latiguillo antes de que el agua
  sea visible. \\[4pt]
Falsos positivos por falta de contexto &
  Sensor Fusion: micrófono MEMS INMP441 (vibración) + sensor de flujo YF-S201 (caudal).
  Si hay ruido pero el grifo está abierto → uso normal. Si hay firma de goteo con
  caudal cero → alerta crítica. \\[4pt]
Dependencia de la nube (latencia, privacidad) &
  Edge AI: la inferencia ocurre íntegramente en el ESP32. Solo el estado de alerta
  (JSON liviano) se transmite por WiFi a la app Flutter. El audio nunca sale del
  dispositivo. Funciona sin internet continuo. \\
\bottomrule
\end{tabularx}

\vspace{6pt}
\begin{mdframed}[linecolor=normalgreen,linewidth=1pt,backgroundcolor=green!3,
  innertopmargin=5pt,innerbottommargin=5pt,innerrightmargin=8pt,innerleftmargin=8pt]
\textbf{Conclusión:} EdgeLeak AI ocupa un nicho no cubierto por ninguno de los
13 referentes: es el primer nodo doméstico de bajo costo que combina
\textbf{prevención acústica + Sensor Fusion + Edge AI} para el punto de mayor
frecuencia de fallo en la cocina residencial (el latiguillo bajo el fregadero).
No reemplaza a las soluciones macro; las miniaturiza donde hacen falta.
\end{mdframed}

% =====================================================================
\section*{\small Referencias}
\vspace{-4pt}
{\footnotesize
\begin{enumerate}[label={[\arabic*]},leftmargin=*,nosep,itemsep=1pt]
\item Islam, M.R. et al. (2023). \textit{An Intelligent IoT and ML-Based Water Leakage Detection System}. IEEE Access, 11. \url{https://ieeexplore.ieee.org/document/10305165}
\item Kang, J. et al. (2023). \textit{ML Model for Leak Detection Using Water Pipeline Vibration Sensor}. Sensors, 23(21). \url{https://doi.org/10.3390/s23218935}
\item Farah, E. \& Shahrour, I. (2024). \textit{Water Leak Detection: A Comprehensive Review}. Water, 16(21). \url{https://doi.org/10.3390/w16212975}
\item Islam, M.R. et al. (2022). \textit{A Review on Water Leakage Detection Technologies}. IEEE Access, 10. \url{https://doi.org/10.1109/ACCESS.2022.3212769}
\item Corado Pérez, L. (2025). \textit{Detección y prevención de fugas en tuberías} [TFM, UOC].
\item Rubiños Jiménez, S.L. et al. (2023). \textit{Sistemas de abastecimiento de agua potable: detección de fugas}. CIDE. \url{https://doi.org/10.33936/978-9942-636-36-2}
\item Bruno, F. et al. (2021). \textit{A Pressure Monitoring System Based on Arduino}. Water, 13(17). \url{https://doi.org/10.3390/w13172321}
\item De Lima Rosado, E.O. (2025). \textit{Fortalecimiento de la seguridad en IoT} [Tesis, UNAD].
\item Dixit, R. et al. (2017). \textit{Water Level and Leakage Detection System}. IRJET, 4(11).
\item Escobar, J. \& López, F. (2016). \textit{Detección de fugas en red de distribución} [Residencia, ITTG].
\item Edge Impulse. (2023). \textit{Anomaly Detection Guide}. \url{https://docs.edgeimpulse.com}
\item Sagat, P. (s.f.). \textit{ESP8266 Nodemcu Leak Sensor} [GitHub]. \url{https://github.com/psagat/ESP8266-Leak-Sensor}
\item Anónimo (s.f.). \textit{Prototipo de detección de fugas domésticas con Arduino} [Trabajo académico].
\end{enumerate}
}

\end{document}
